% use \chaptermark{}
% to alter or adjust the chapter heading in the running head


%%%%%%%%%%%%%%%%%%%%%%%%%%%%%%%%%%%%%%%%%%%%%%%%%%%%%%%%%%
%%%
%%%  The text below comes from (Feinsilver)
%%%  Trace Identities and Integration of Parts for Zeons
%%%
%%%%%%%%%%%%%%%%%%%%%%%%%%%%%%%%%%%%%%%%%%%%%%%%%%%%%%%%%%

\section{Zeon algebra and structure of the kernel}
\label{sec:7.1}

\subsection{Introduction}
\label{sec:7.1.1}

The notation $\diag{v}$ denotes the diagonal matrix with entries $v_i$ of the vector $\bf v$.

The \textsl{theta symbol\/} for pairs of single indices
$$
  \theta_{ij}=\begin{cases} 
	1, & \text{if } i<j\cr  
	0, & \text{otherwise} \end{cases}
$$

We indicate the component of a vector by underlining and then subscripting, e.g. $\underline{Tv}_{\,i}$
is the $i^{\rm th}$ component of the result of the matrix $T$ acting on the vector $v$.

$\V$ is the vector space isomorphic to $\BDQ^n$. Denote by ${\rm Sym}_n(\V)$ the space of
symmetric matrices acting on $\V$.

Denote by $\Z$ the algebra of \textsl{zeons}, generated by $n$ commuting variables $e_i$ such that
$e_i^2=0$, $1\le i \le n$. Then we identify $\V$ with the $\BDQ$-span of the zeons and generate
an algebra, $\Z$. The $k^{\rm th}$ zeon tensor power of $\V$ is denoted $\spow{\V}{k}$. 
Effectively $\spow{\V}{k}$ is the space of symmetric tensors with no repeated factors.
Elements of $\spow{\V}{k}$ are indexed by $k$-tuples of positive integers all distinct. 
I.e., the space has dimension $\binom{n}{k}$. We will denote multi-indices by 
roman capital letters $\rI=(i_1,i_2,\ldots,i_k)$, $\rJ$, $\rK$, etc. 
A linear operator $W$ on $\V$ extends to the operator $\spow{W}{k}$ acting on $\spow{\V}{k}$.
The \textsl{zeon Fock space\/}, 
$$ 
  {\mathcal F}=\BDQ\oplus (\bigoplus_{k\ge1}\spow{\V}{k})
$$
is finite dimensional, like the exterior algebra, and the \textsl{second quantization\/} of $W$ is the extension of 
$W$ to $\mathcal F$. Of course, one could work over a different field such as $\BDC$ for quantum mechanics with zeons.

The $\rI\rJ$ component of $\spow{W}{k}$ is the permanent of the corresponding submatrix of $W$ with
rows indexed by $\rI$ and columns by $\rJ$. For dimension 2, we will use indexes $ij$ instead
of $(i_1,i_2)$, as usual.

For given $n$, $X$, $Y$, etc., are vectors in $\spow{\V}{2}\approx\BDQ^{\binom{n}{2}}$.
The mapping \hfb ${\rm Mat}\colon\,\spow{\V}{2}\to {\rm Sym}_n(\V)$ is a linear embedding taking
the vector $X=(x_{ij})$ to the symmetric matrix $\hat X$ with off-diagonal components 
$\hat X_{ij}=x_{ij}=\hat X_{ji}$, for $i<j$ and zeros on the diagonal, $\hat X_{ii}=0$, $1\le i\le n$.
We will use the explicit notation ${\rm Mat}(X)$ as needed for clarity.

We equip $\spow{\V}{2}$ with the inner product 
$$
  \avg X,Y.=X_{\lambda\mu} Y_{\lambda\mu}=\tr \hat X\hat Y\ .
$$


%%% Basic Identities
\subsection{Basic Identities}
\label{sec:7.1.2}

Multiplying $X$ with $u$, we observe that
%
\begin{equation}\label{eqtrace1}
  X\dg u.=\frac12\,\tr \hat X J=u\dg  X.
\end{equation}
%
Observe also that if $D$ is diagonal, then $\tr D=\tr DJ$.

\begin{proposition}\label{propBR}   	\label{prop:6.6}
(\textsl{Basic Relations}) We have
\begin{enumerate}
\item[{\normalfont 1.}] $\Mat(X\spow{A}{2}) = \symm{A}{\hat X}-D^+$, where $D^+$ is a diagonal matrix satisfying $\vstrut\tr D^+=\tr \symm{A}{\hat X}$.

\item[{\normalfont 2.}] $\Mat(\spow{A}{2}\dg  X.) = \sym{A}{\hat X}-D^-$, where $D^-$ is a diagonal matrix satisfying $\vstrut\tr D^-=\tr \sym{A}{\hat X}$.

\item[{\normalfont 3.}] If $A$ and $X$ have nonnegative entries, then $D^+$ and $D^-$ have nonnegative entries. In particular, in that case,
vanishing trace for $D^{\pm}$ implies vanishing of the corresponding matrix.
\end{enumerate}
\end{proposition}

\begin{proof} The components of $X\spow{A}{2}$ are
%
\begin{eqnarray*}
\underline{X\spow{A}{2}}_{\,ij}  
  &=& \theta_{ij}\theta_{\lambda\mu}
	(x_{\lambda\mu}A_{\lambda i}A_{\mu j}
	+x_{\lambda\mu}A_{\mu i}A_{\lambda j}) \\
  &=& \theta_{ij} x_{\lambda\mu} A_{\lambda i}A_{\mu j}
	=\theta_{ij}\underline{\symm{A}{\hat X}}_{\,ij}
\end{eqnarray*}
%
noting that the diagonal terms of $\hat X$ vanish anyway. And $\symm{A}{\hat X}$ will be symmetric if $\hat X$ is.
Since the left-hand side has zero diagonal entries, we can remove the theta symbol and compensate by subtracting off
the diagonal, call it $D^+$. 
Taking traces yields \#1. And \#2 follows similarly.
\end{proof}

\begin{remark}   	\label{rem:7.1}
Observe that $D^+$ and $D^-$ may be explicitly given by
%
\begin{eqnarray*}
  D^+_{\,ii}  &=& x_{\lambda\mu}A_{\lambda i}A_{\mu i} \\
  D^-_{\,ii}  &=& x_{\lambda\mu}A_{i\lambda }A_{i\mu} 
\end{eqnarray*}
%
where for $D^+$ the $A$ elements are taken within a given column, while for $D^-$, the $A$ elements are in a given row.
\end{remark}

\begin{proposition}\label{propXX}  	\label{prop:7.2} %old 6.7
Let $X$ and $A$ be nonnegative. Then 
%
\begin{eqnarray*}
  \hat X = \symm{A}{\hat X}  &\Rightarrow&  X\spow{A}{2} = X\\
  \hat X = \sym{A}{\hat X}  &\Rightarrow&  \spow{A}{2}\dg  X. = \dg  X.
\end{eqnarray*}
%
\end{proposition}

\begin{proof} Start with
$$
  D^+= \symm{A}{\hat X}-\Mat(X\spow{A}{2})
$$
with $D^+$ vanishing if $\tr D^+=\tr \symm{A}{\hat X}$ vanishes.
If $\hat X = \symm{A}{\hat X}$, then since $\hat X$ has vanishing trace, $\tr \symm{A}{\hat X}=0$. 
Thus, $\tr D^+=0$ and $\hat X = \symm{A}{\hat X}=\Mat(X\spow{A}{2})$. 
The second implication follows similarly. 
\end{proof}

%%% Trace identities
\subsection{Trace Identities}
\label{sec:7.1.3}

Using equation (\ref{eqtrace1}), we will find some identities for these quantities.

\begin{proposition}\label{propTR}  \label{prop:7.3}
We have 
\begin{enumerate}
\item[{\normalfont 1.}] $\displaystyle X\spow{A}{2}\dg  u. = \half \,\tr (\hat X A(J-I)A^*)$ .

\item[{\normalfont 2.}] $\displaystyle u\spow{A}{2}\dg  X. = \half \,\tr (\hat X A^*(J-I)A)$ .

\item[{\normalfont 3.}] If $A$ is stochastic, then \ $X\spow{A}{2}\dg  u. = \half \,\tr (\hat X (J-AA^*))$.
\end{enumerate}
\end{proposition}

\begin{proof} We have, using equation (\ref{eqtrace1}) and Basic Relation 1,
%
\begin{eqnarray*}
  X\spow{A}{2}\dg  u.  
 &=& \half \,\tr \Mat(X\spow{A}{2})J\\
 &=& \half\,\tr (\symm{A}{\hat X}J-D^+J) \\
 &=& \half\,\tr (\symm{A}{\hat X}J-\symm{A}{\hat X})
 \end{eqnarray*}
%
and rearranging terms inside the trace yields \#1. 
Then \#3 follows since $A$ stochastic implies $AJ=J=JA^*$.
And \#2 follows similarly, using the second Basic Relation in the equation
$\displaystyle u\spow{A}{2}\dg  X. = \half \,\tr \Mat(\spow{A}{2}\dg  X.)J$.\qed
\end{proof}

Using  equation (\ref{eqtrace1}) directly for $X$, we have 
%
\begin{eqnarray}
X(I-\spow{A}{2})\dg  u.  &=&  \half \,\tr (\hat X (J-AJA^*+AA^*)) \label{eqbp1}\\
u(I-\spow{A}{2})\dg  X.  &=&  \half \,\tr (\hat X (J-A^*JA+A^*A))\label{eqbp2}
\end{eqnarray}
%


%%% stochastic case
\subsubsection{Stochastic case}
\label{sec:7.1.3.1}

For stochastic $A$, equation (\ref{eqbp1}) yields
\begin{equation}\label{eqparts}
X(I-\spow{A}{2})\dg  u. = \half \,\tr A^*\hat XA
\end{equation}

The above formula is ``integration-by-parts for zeons".

We denote an invariant distribution for $A$ by $\pi$, i.e. $\pi A=\pi=(p_1,\ldots,p_n)$.
We now have a converse result to Proposition 2. 

\begin{proposition}\label{propST}  	\label{prop:7.4}  %old 6.9
Let  $A$ be stochastic and $X$ nonnegative.
\begin{enumerate}
\item[{\normalfont 1.}] $X\spow{A}{2} = X \Rightarrow  \hat X = \symm{A}{\hat X}$.

\item[{\normalfont 2.}] If $A$ has a strictly positive invariant distribution, $\pi$, then \\
$\spow{A}{2}\dg  X. = \dg  X.   \Rightarrow  \hat X = \sym{A}{\hat X}\vstrut$.
\end{enumerate}
\end{proposition}

\begin{proof} If $X\spow{A}{2} = X$, then Proposition \ref{propBR}, eq. \eqref{eq:N}, reads
$$
  \hat X = \symm{A}{\hat X}-D^+
$$
Multiply by $\dg  u.$ on the right and by $u$ on the left, using  $A\dg  u.=\dg  u.$ and $uA^*=u$ to get
$$ 
  u\hat X\dg  u. =u\hat X\dg  u.-uD^+\dg  u.
$$
So $uD^+\dg  u.=0$ and taking traces yields $\tr(D^+J)=0$, thus $D^+=0$.

If $\spow{A}{2}\dg  X. = \dg  X. $, then Proposition \ref{propBR}, eq. \eqref{eq:M}, reads
$$
  \hat  X = \sym{A}{\hat X}-D^-
$$
Multiply by $\dg  \pi.$ on the right and by $\pi$ on the left, using $\pi A=\pi$ and $A^*\dg \pi.=\dg \pi.$ to get
$$
  \pi\hat X\dg  \pi. =\pi\hat X\dg  \pi.-\pi D^-\dg  \pi.
$$
So $\pi D^-\dg  \pi.=0$ and taking traces yields $\tr(D^-\dg  \pi.\pi)=0$. I.e., $\displaystyle \sum_i D^-_{ii}p_i^2=0$, so
$D^-=0$. \qed
\end{proof}

With $A$ stochastic, irreducible and aperiodic, $\displaystyle \lim_{n\to\infty}A^n = \Omega=\dg  u.\pi$, satisfying
$\Omega^2=\Omega=A\Omega=\Omega A$. Observe that
%
\begin{eqnarray*}
  \Omega\Omega^* &=& \dg  u.\pi\dg  \pi.u =(\sum p_i^2)\, J\\
  \Omega^*\Omega &=& \dg  \pi.u\dg  u.\pi = n\,\dg  \pi.\pi
\end{eqnarray*}
%
both of these having all positive entries.  Thus

\begin{proposition}\label{propIRPER} 	\label{prop:7.5}  %old 6.10 
Let $A$ be stochastic, irreducible and aperiodic. Then, for nonnegative $X$, either of 
$\hat X=\sym{A}{\hat X}$ or $\hat X=\symm{A}{\hat X}$ implies $X=0$.
\end{proposition}

\begin{proof} If $\hat X=\sym{A}{\hat X}$, then, inductively for positive integers $k$,
$$ 
  \hat X=A^k\hat X (A^*)^k
$$
Letting $k\to\infty$, we have $\hat X=\Omega \hat X \Omega^*$. Taking traces yields $\tr \hat X \Omega^*\Omega=0$. 
Since $\hat X$ is nonnegative and $\Omega^*\Omega$ is a matrix with positive entries, we have $X=0$. 
The proof for $\hat X=\symm{A}{\hat X}$ is similar.  \qed
\end{proof}

Now we have 

\begin{theorem}\label{thmIRPER}
Let $\spow{A}{2}\ne0$. 
Then\\ 
$A$ is aperiodic and irreducible if and only if $\det(I-\spow{A}{2})\ne 0$.
\end{theorem}

\begin{remark}		\label{rem:7.2}
Note that the hypothesis is saying that $A$ is not the identity matrix, i.e., that the chain moves at least one state to some other one.
\end{remark}

\begin{proof}
If $A$ is reducible, let $\{C_0,\ldots,C_{s-1}\}$ denote the closed classes partitioning the states.
Define $X$ by $x_{ij}=0$ if there is a $k$ such that $i,j\in C_k$, $1$ otherwise, i.e.,
$$ 
  X  = 1-\sum_k \chi_{C_k\times C_k}
$$
We will show that $\spow{A}{2}\dg  X.=\dg  X.$. 
\begin{description}
\item[Step 1.] First check that the diagonal vanishes, i.e., $A_{i\lambda}x_{\lambda\mu}A_{i\mu}=0$. 
By assumption, if $i\in C_k$, then
$\lambda,\mu\in C_k$. 
But then $x_{\lambda\mu}=0$. \\

\item[Step 2.] For $i\ne j$, $\underline{\spow{A}{2}\dg  X.}_{\,ij}=A_{i\lambda}x_{\lambda\mu}A_{j\mu}$. 
If $i\in C_k$, then the sum is over
$\lambda\in C_k$. If $j\in C_k$, then the sum is over $\mu\in C_k$ as well, $x_{\lambda\mu}=0$, and the result is zero.
If $j\in C_{k'}$, $k'\ne k$, then $x_{ij}=1$. 
Writing out the sum, we have
$$ 
  \sum_{l\in C_k, m\in C_{k'}} A_{il}x_{lm}A_{jm}
$$
For fixed $m$, $x_{lm}=1$, and the sum over $l\in C_k$ gives all choices that $i$ could map to, while for $l\notin C_k$, $A_{il}=0$, so 
$\displaystyle  \sum_{l\in C_k} A_{il}=1$. 
Similarly, summing over $m$ gives 1. 
Hence, we get $1$ in case $x_{ij}\ne0$ and
$0$ otherwise, as required. 
\end{description}

If $A$ is periodic, with period $p$, say, let $\{C_0,\ldots,C_{p-1}\}$ denote the classes partitioning the states such that 
$i\in C_k$ implies $iA\in C_{k+1}$, with indices modulo $p$.
Define $X$ by $x_{ij}=1$ if there exists $k$ such that $i\in C_k$ and $j\in C_{k+1}$, $0$ otherwise.  We check that 
$\spow{A}{2}\dg  X.=\dg  X.$.  The steps are similar to that for the case of reducibility:\\

Step 1. First check that the diagonal vanishes, i.e., $A_{i\lambda}x_{\lambda\mu}A_{i\mu}=0$. By periodicity, if $i\in C_{k-1}$, then
$\lambda\in C_k$ and hence $\mu\in C_{k+1}$ by definition of $X$. But then $A_{i\mu}$ vanishes since $\mu\in C_{k+1}$ requires
$i\in C_k$.\\

Step 2. For $i\ne j$, $\underline{\spow{A}{2}\dg  X.}_{\,ij}=A_{i\lambda}x_{\lambda\mu}A_{j\mu}$. If $i\in C_{k-1}$, then
$\lambda\in C_k$, $\mu\in C_{k+1}$ requires $j\in C_k$ so $x_{ij}\ne 0$. Now, writing out the sum, we have
$$ 
  \sum_{l\in C_k, m\in C_{k+1}} A_{il}x_{lm}A_{jm}
$$
For fixed $m$, $x_{lm}=1$, and the sum over $l$ gives all choices that $i$ could map to, i.e.,
$\displaystyle  \sum_{l\in C_k} A_{il}=1$. And summing over $m$, gives 1 as well. Hence, we get $1$ in case $x_{ij}\ne0$ and
$0$ otherwise, as required.\\

Suppose $A$ is aperiodic and irreducible.  By Perron-Frobenius the top nonnegative eigenvalue has a nonnegative eigenvector, $X$.
Proposition \ref{propST} says that a solution to $X\spow{A}{2} =X$ satisfies $\hat X=\symm{A}{\hat X}$. 
Then Proposition \ref{propIRPER} implies that $X=0$.
\end{proof}

%%% CSS

\subsection{Kernel structure}
\label{sec:7.1.4}

Let $\displaystyle A_2=\spow{A}{2}+\spow{\Delta}{2}=\half\,(\spow{R}{2}+\spow{B}{2})$.

For $\Delta$ we have, from Proposition \ref{propTR},
$$
  X\spow{\Delta}{2}\dg  u. 
  = \half \,\tr (\hat X (\Delta J\Delta^*-\Delta \Delta^*))
$$
which becomes
$$
  X\spow{\Delta}{2}\dg  u. = -\half \,\tr \Delta^*\hat X\Delta
$$
since $\Delta J=0$. And Proposition \ref{propTR}, \#2, yields
$$
  u\spow{\Delta}{2}\dg  X. 
	= \half \,\tr (\hat X (\Delta^*(J-I)\Delta)\ .
$$


\subsubsection{$\mathbf M$ and $\mathbf N$}
\label{7.1.4.1}

Note that $\Omega=\avg K.$, the average element of the kernel. As well, $\Omega=\beta_\lambda\rho(k_\lambda)$,
where $k_i$ runs through any row, so that each range is represented once.

Let $\omega$ be the diagonal matrix with $\pi$ on the diagonal, i.e. $\omega=\diag{\pi}$. 

\begin{proposition} \label{propKstarK} 	\label{prop:7.6}  %old 6.11
$\avg K^*K.=n\,\omega$.
\end{proposition}
\begin{proof}
For \#1, start by observing that $\diag{v}=\diag{\diag{v}{\rm \dg  u.}}$. Thus, since $k^*k$ is diagonal for every $k$, we check 
$\avg K^*K{\dg  u.}.=\avg K^*{\dg  u.}.=\Omega^*\dg  u.=n\,\dg  \pi.\ .$
\end{proof}

If  $\Omega_2$ is the limiting matrix of $A_2$, define $u_2=\Omega_2\dg  u.$ and $\pi_2=u\Omega_2$.
Then we have the following relations.

\begin{proposition}\label{propMN}  	\label{prop:7.7}  %old 6.12
For any $X$,
$$
  X\Omega_2\dg  u.=X\dg  u_2.=\avg X,u_2.
	=\half\,\tr \hat X\hat u_2 = \half\,\tr \hat XN
$$
where $N=\avg J-KK^*.$, taken over the kernel. \\
And
$$
  u\Omega_2\dg  X.=\pi_2\dg  X.
	=\avg \pi_2,X.=\half\,\tr \hat \pi_2\hat X 
	= \half\,\tr \hat X\tilde M
$$
where $\tilde M=\avg K^*(J-I)K.$, taken over the kernel.
\end{proposition}

\begin{proof} Using the homomorphism property of the map $\spow{\cdot}{2}$, we have $\Omega_2=\avg \spow{K}{2}.$ 
taken over the kernel. 
Thus, using Proposition \ref{propTR} for stochastic matrices, 
$$
  X\Omega_2\dg  u.=\avg X\spow{K}{2}{\dg  u.}.
	=\half\, \tr \avg \hat X(J-KK^*).= \half\,\tr \hat XN
$$
Similarly,
$$
  u\Omega_2\dg  X.=\avg u\spow{K}{2}{\dg  X.}.
	=\half\, \tr \avg \hat XK^*(J-I)K).
	= \half\,\tr \hat X\tilde M\ . \qed
$$ 
\end{proof}

Let $N_0=\avg KK^*.  = \avg EE^*.$, where $K$ denotes a random element of $\cK$, and $E$ denotes a random idempotent,
element of $\E(\cK)$. 
Define $N=J-N_0$. 
Since $\avg J.=J$, we set $u_2=\Omega_2\dg  u.$, satisfying $A_2u_2=u_2$, and identify $N$ as $\avg J-KK^*.=\Mat(u_2)$.
Note that $N=\alpha_\mu e_\mu e_\mu^*$ where the $e_i$ are (all of) the idempotents with a given range, i.e., they can be chosen as the idempotents from a single column in the kernel.

Let $M_0=\avg K^*JK.$. Define $\tau=\tr M_0$ and  $M=M_0-\tau\omega$.  Since $M$ and $\tilde M$ agree on off-diagonal 
elements, we set $\pi_2=u\Omega_2$, satisfying $\pi_2A_2=\pi_2$, and identify $M$ as $\Mat(\pi_2)$.

Note that $M_0$, $N_0$, $M$, $N$ as well as $J$ are all symmetric matrices.

Define the mappings, $F,G\colon {\rm Sym}_n(\V)\to{\rm Sym}_n(\V)$ by
$$ 
  F(Y)=\sym{A}{Y}+\sym{\Delta}{Y} \qquad\hbox{ and }\qquad
	G(Y)=\symm{A}{Y}+\symm{\Delta}{Y}\ .
$$
Note that
$$ 
  F(Y)=\half\,(\sym{R}{Y}+\sym{B}{Y}) \qquad\hbox{ and }\qquad
	G(Y)=\half\,(\symm{R}{Y}+\symm{B}{Y})\ .
$$
First, some observations

\begin{proposition}   	\label{prop:7.8}   %old 6.13
For any vector $v$, \rm $G(\diag{v}) = \diag{vA}$ .
\end{proposition}

\begin{proof} Let $Y=G(\diag{v})$. 
Then $2\,Y_{ij}=R_{\lambda i}v_{\lambda}R_{\lambda j}+B_{\lambda i}v_{\lambda}B_{\lambda j}$.
A term like $R_{li}$ means that $R$ maps $l$ to $i$. 
Since $R$ is a function, $l$ can only map to both $i$ and $j$ if $i=j$.
So $Y$ is diagonal. If $i=j$, then we get, using $R_{ij}^2=R_{ij}$, $B_{ij}^2=B_{ij}$,
$$
  2\,Y_{ii}=R_{\lambda i}v_{\lambda}+B_{\lambda i}v_{\lambda}
	=2\,v_\lambda A_{\lambda i}
$$
as required.
\end{proof}

For example, taking $v=u$, we see that 
$$
  G(I)=I \ \hbox{if and only if}\  A \ 
	\hbox{is doubly stochastic.}
$$

\begin{proposition}  	\label{prop:7.9}  %old 6.14
$F(J)=J$ and $G(\omega)=\omega$.
\end{proposition}

\begin{proof}
First, we have $F(J)=\sym{A}{J}+\sym{\Delta}{J}$, with $AJ=J=JA^*$ and $\Delta J=0$, this reduces to $J$.
Second, from the previous Proposition, we have $G(\omega)=G(\diag{\pi}))=\diag{\pi A}=\diag{\pi}=\omega$.
\end{proof}

Note that equation (\ref{eqparts}) applied to $R$ and $B$ and then averaged yields
%
\begin{equation} \label{eqPTS}
  X(I-A_2)\dg  u. =\half\,\tr G(\hat X)
\end{equation}
%

The next theorem shows the invariance of $N$ and $M$ under $F$ and $G$ respectively as well as the general relation between
fixed points of $A_2$ and fixed points of $F$ and $G$. 
Note that we are not restricting {\it a priori\/} to nonnegative solutions except in case 2 of the following theorem.

\begin{theorem}   \label{thm:7.2}  \hfill \\[-8pt]
\begin{enumerate}
\item[{\normalfont 1.}] $F(\hat X)=\hat X$ if and only if $A_2\dg  X.=\dg  X.$.

\item[{\normalfont 2.}] For nonnegative solutions $X$, $G(\hat X)=\hat X$ if and only if $XA_2=X$.
\end{enumerate}
\end{theorem}

\begin{proof}
First, from Proposition \ref{propBR}, applying the Basic Relations for $R$ and $B$ and averaging, we get
%
\begin{eqnarray}
\Mat(XA_2)  &=& G(\hat X)-\half\,(D_R^++D_B^+) \label{eqG}\\
\Mat(A_2\dg  X.)  &=&  F(\hat X)-\half\,(D_R^-+D_B^-)
\end{eqnarray}
%
where the subscripts on the $D$'s indicate the diagonals corresponding to $R$ and $B$ respectively.
Recalling the remark following Proposition \ref{propBR}, note that for $D_R^-$, say, we have terms like
$R_{i\lambda}x_{\lambda\mu}R_{i\mu}$. Since $R$ is a function, we must have $\lambda=\mu$, but then $x_{\lambda\mu}=0$.
In other words, $D_R^-$ and $D_B^-$ vanish so that, in fact,  $\Mat(A_2\dg  X.)  =  F(\hat X)$. And \#1 follows immediately.

For \#2,  if $G(\hat X)=\hat X$, then the trace of the diagonal terms in equation (\ref{eqG}) vanish, so under the assumption of
nonnegativity, they must vanish. And replacing $G(\hat X)$ by $\hat X$, we have $XA_2=X$.

Assume $XA_2=X$. Taking traces in equation \eqref{eqG}, the trace of the left-hand side vanishes and the trace of the diagonal
terms equals $\tr G(\hat X)$.  Now, use the integration-by-parts formula for $A_2$, equation \eqref{eqPTS}. The left-hand
side vanishes, hence $\tr G(\hat X)=0$, implying the vanishing of the terms $D_R^+$ and $D_B^+$ as well. Then equation
\eqref{eqG} reads $\Mat(X)=G(\hat X)$ as required. \qed
\end{proof}

From this Theorem we have \\

With $\displaystyle A_2=\spow{A}{2}+\spow{\Delta}{2}=\half(\spow{R}{2}+\spow{B}{2})$, 
we have $\Omega_2$ as the Ces\`aro limit of $A_2^k$.
Then with $u$ here denoting the all-ones vector of dimension $\binom{n}{2}$,
$u_2=\Omega_2\dg  u.$ and $\pi_2=u\Omega_2$. Then
\begin{enumerate}
\item[{\normalfont 1.}] $N=\avg J-KK^*.=\Mat(u_2)$ satisfies $\sym{A}{N}+\sym{\Delta}{N}=N$. It is a nonnegative solution with $\tr N=0$.
$N_0=\avg KK^*.$ and $J=\dg  u.u$ are also nonnegative solutions, with trace equal to $n$.\\

\item[{\normalfont 2.}] $M=\avg K^*(J-(\tau/n) I)K. =\Mat(\pi_2)$ satisfies $\symm{A}{M}+\symm{\Delta}{M}=M$. It is a nonnegative solution with $\tr M=0$.
$M_0=\avg K^*JK.$ and $\omega=\diag{\pi}$ are also nonnegative solutions, with trace equal to $\tau$ and $1$ respectively.
\end{enumerate}

Note that in Proposition \ref{propMN} we get $\tilde M=\avg K^*(J-I)K.$. However, $\tr \tilde M$ will not be zero. To cancel off the diagonal of $M_0$, we need to subtract off a multiple of $\omega$. 
From Proposition \ref{propKstarK}, we see that an appropriate multiple of $K^*K$ will work.


\subsubsection{Determinants and traces}
\label{sec:7.1.4.2}

Let $X=e_1\,{\rm adj}(I-\spow{A}{2})$ in equation \eqref{eqparts}. Since $X(I-\spow{A}{2})\dg  u.=\det(I-\spow{A}{2})$, we get
$$ 
  \det(I-\spow{A}{2})=\half\,\tr \symm{A}{\hat X}
$$

Similarly, let $X=e_1\,{\rm adj}(I-A_2)$ in equation \eqref{eqPTS}. 
We get
$$ 
  \det(I-A_2)=\half\,\tr G(\hat X)
	=\half\,\tr(\hat X(AA^*+ \Delta\Delta^*))
$$
Note that if $\det(I-A_2)=0$, then, if nonvanishing, $X=e_1\,{\rm adj}(I-A_2)$ will be a nontrivial solution to
$XA_2=X$. 
That is, if the solution space is one-dimensional, this is a multiple of $\pi_2$ and we scale to $\hat X=M$.
Then the vanishing of the trace will be equivalent to the vanishing of the determinant, i.e., the existence of solutions to $XA_2=X$.
